% FAC-EN-U0063
% Source: fac_body.tex lines 2406--2430.
% Authority: PDF physical pp. 63--64; printed pp. 258--259
% (Chapter III, no. 65 complete).

\phantomsection\label{fac:III:3:no:65}
\textbf{65. Applications.}

\phantomsection\label{fac:III:3:no:65:prop:5}
\textsc{Proposition 5.} \textit{If $M$ is a graded $S$-module
\hyperref[fac:III:2:no:56:def:TF]{satisfying condition}}
\hyperref[fac:III:2:no:56:def:TF]{(TF)}\textit{, the homomorphism
$\alpha:M\to\Gamma(\mathcal{A}(M))$, defined in
\hyperref[fac:III:2:no:59:def:alpha]{no.~59}, is
$\mathcal{C}$-bijective.}

We must show that
$\alpha:M_n\to\Gamma(X,\mathcal{A}(M(n)))$ is bijective for all
sufficiently large $n$. By
\hyperref[fac:III:3:no:64:prop:4]{Proposition~4}, however,
$\Gamma(X,\mathcal{A}(M(n)))$ is identified with $H^0(M(n))$.
\hyperref[fac:III:3:no:65:prop:5]{The proposition} therefore follows from
\hyperref[fac:III:3:no:63:prop:3:item:c]{Proposition~3(c)}, taking into
account that, under the preceding identification, the homomorphism
$\alpha$ is transformed into the homomorphism defined at the beginning of
\hyperref[fac:III:3:no:62:def:alpha]{no.~62}, which is also denoted by
$\alpha$.

\phantomsection\label{fac:III:3:no:65:prop:6}
\textsc{Proposition 6.} \textit{Let $\mathcal{F}$ be a coherent algebraic
sheaf on $X$. The graded $S$-module $\Gamma(\mathcal{F})$
\hyperref[fac:III:2:no:56:def:TF]{satisfies condition}}
\hyperref[fac:III:2:no:56:def:TF]{(TF)}\textit{, and the homomorphism
$\beta:\mathcal{A}(\Gamma(\mathcal{F}))\to\mathcal{F}$, defined in
\hyperref[fac:III:2:no:59:def:beta]{no.~59}, is bijective.}

By \hyperref[fac:III:2:no:60:theorem:2]{Theorem~2 of no.~60}, we may
suppose that $\mathcal{F}=\mathcal{A}(M)$, where $M$ is a module
\hyperref[fac:III:2:no:56:def:TF]{satisfying (TF)}. By
\hyperref[fac:III:3:no:65:prop:5]{the preceding proposition},
$\alpha:M\to\Gamma(\mathcal{A}(M))$ is
\hyperref[fac:III:2:no:56:def:C-morphism]{$\mathcal{C}$-bijective}. Since
$M$ \hyperref[fac:III:2:no:56:def:TF]{satisfies (TF)}, it follows that
$\Gamma(\mathcal{A}(M))$ does also. Applying
\hyperref[fac:III:2:no:58:prop:6]{Proposition~6 of no.~58}, we see that
\[
\alpha:\mathcal{A}(M)\longrightarrow
\mathcal{A}(\Gamma(\mathcal{A}(M)))
\]
is bijective. Since the composite
\phantomsection\label{fac:III:3:no:65:formula:alpha-beta-composite}
\[
\mathcal{A}(M)\xrightarrow{\ \alpha\ }
\mathcal{A}(\Gamma(\mathcal{A}(M)))
\xrightarrow{\ \beta\ }\mathcal{A}(M)
\]
is the identity
(\hyperref[fac:III:2:no:59:prop:7]{no.~59, Proposition~7}), it follows that
$\beta$ is bijective, which proves
\hyperref[fac:III:3:no:65:prop:6]{the proposition}.

\phantomsection\label{fac:III:3:no:65:prop:7}
\textsc{Proposition 7.} \textit{Let $\mathcal{F}$ be a coherent algebraic
sheaf on $X$. The groups $H^q(X,\mathcal{F})$ are finite-dimensional
vector spaces over $K$ for every $q\geq0$, and
$H^q(X,\mathcal{F}(n))=0$ for $q>0$ and all sufficiently large $n$.}

As above, we may suppose that $\mathcal{F}=\mathcal{A}(M)$, where $M$ is a
module \hyperref[fac:III:2:no:56:def:TF]{satisfying (TF)}.
\hyperref[fac:III:3:no:65:prop:7]{The proposition} then follows from
\hyperref[fac:III:3:no:63:prop:3]{Proposition~3} and the
\hyperref[fac:III:3:no:64:cor]{corollary to Proposition~4}.

\phantomsection\label{fac:III:3:no:65:prop:8}
\textsc{Proposition 8.} \textit{We have
$H^q(X,\mathcal{O}(n))=0$ for $0<q<r$, and
$H^r(X,\mathcal{O}(n))$ is a vector space of dimension
$\binom{-n-1}{r}$ over $K$, with basis given by the cohomology classes of
the alternating cocycles of $\mathfrak{U}$}
\phantomsection\label{fac:III:3:no:65:formula:top-basis}
\[
f_{01\cdots r}=\frac{1}{t_0^{\beta_0}\cdots t_r^{\beta_r}},
\qquad \textit{where}\qquad
\beta_i>0
\quad\textit{and}\quad
\sum_{i=0}^{i=r}\beta_i=-n.
\]
We have $\mathcal{O}(n)=\mathcal{A}(S(n))$, whence
$H^q(X,\mathcal{O}(n))=H^q(S(n))$ by the
\hyperref[fac:III:3:no:64:cor]{corollary to Proposition~4}.
\hyperref[fac:III:3:no:65:prop:8]{The proposition} now follows immediately
from this and the corollaries to
\hyperref[fac:III:3:no:62:prop:2]{Proposition~2}.

In particular, $H^r(X,\mathcal{O}(-r-1))$ is a one-dimensional vector
space over $K$, with basis given by the cohomology class of the cocycle
\phantomsection\label{fac:III:3:no:65:formula:canonical-top-class}
\[
f_{01\cdots r}=\frac{1}{t_0\cdots t_r}.
\]
